\section{Introduction}
\label{sec:intro}

In-context learning (ICL) allows foundation models, especially large language models (LLMs), to adapt to new tasks from only a few input-output demonstrations, without updating model parameters~\cite{min2022rethinking, dong2022survey, highmore2024context, wang2023large, sia2024does, li2025large}.
Recent progress in large vision-language models (VLMs) has extended this paradigm from language to multimodal settings, giving rise to visual in-context learning (VICL), where a model infers an image transformation directly from a small number of visual examples~\cite{bar2022visual, wang2023painter, wang2023incontext, xprompt2024universal, zhou2024visual, ma2024vision}.
Such a capability is particularly appealing for low-level vision, where one would ideally like a single frozen model to handle diverse image restoration, removal, and generation tasks by simply conditioning on demonstrations rather than retraining for each task.

However, existing VICL methods largely rely on a restrictive \emph{same-task} assumption: the demonstration pair and the query are usually drawn from the same task. In practice, this assumption is often violated. For example, a user may only have a \emph{deraining} example while the query requires \emph{dehazing}, or may provide a reflection-removal pair while the true target is shadow removal. This mismatch makes VICL fundamentally ambiguous: should the model imitate the demonstrated transformation, or infer a different one from the query itself? At the same time, many low-level tasks are not isolated~\cite{zamir2018taskonomy, pal2019tmt, achille2019task2vec, bao2019information}. They share latent relationships in terms of visual effects, such as artifact suppression, contrast recovery, color adjustment, or illumination correction. Importantly, these relationships are often easier to express in language than in raw pixels, suggesting that language may serve as a natural bridge for transferring knowledge across mismatched visual tasks~\cite{trigka2025comprehensive, shu2024visual, brooks2023instructpix2pix, potlapalli2023promptir, conde2024instructir, ke2023neural}.

These observations motivate a new question: \emph{Can VLMs perform visual in-context learning when the demonstration and the query come from different tasks?} We study this setting as \textbf{cross-task visual in-context learning}. Compared with standard VICL, cross-task VICL requires more than pattern imitation. The model must abstract the visual change encoded in the context, reason about the discrepancy between the demonstrated task and the query task, and then produce the correct target transformation without being explicitly told the task name. A fixed prompt is often insufficient for this setting, because it provides only a generic instruction while leaving the task difference unresolved. What is needed is a mechanism that converts mismatched visual demonstrations into task-aware guidance that remains compatible with frozen VLM generators.

To this end, we propose \textbf{T2T-VICL}, a collaborative teacher-student prompting framework for cross-task VICL. Our key idea is to translate the relationship between two tasks into \emph{implicit task-difference prompts} that describe visual changes in a task-agnostic manner. A large teacher VLM first generates structured implicit descriptions from cross-task image tuples, and we curate these outputs into a benchmark of cross-task text-image relations. This capability is then distilled into a lightweight student VLM, which produces content-dependent prompts from a task-$A$ demonstration pair and a task-$B$ query. The generated prompt is fed to a frozen image-editing VLM, while an automatic score-based inference strategy, combining attribution-enhanced prompting, VLM-based visual instruction scoring~\cite{ku2023viescore}, and image quality assessment (IQA) metrics~\cite{zhang2012comprehensive}, is used to rank candidate outputs. In this way, T2T-VICL provides a practical framework for probing and improving cross-task VICL without updating the final generator model.

Our contributions are three-fold:
\begin{itemize}[itemsep=-2pt, topsep=1pt]
    \item  We formulate \textbf{cross-task VICL} for low-level vision and introduce a benchmark with implicit cross-task descriptions, enabling systematic study of how VLMs generalize across mismatched visual demonstrations. 
    \item  We propose T2T-VICL, a \textbf{VLM $\rightarrow$ sVLM $\rightarrow$ VLM} teacher-student pipeline that distills implicit task-relation knowledge from a large model into a lightweight prompt generator for frozen image-editing VLMs. 
    \item  We develop a \textbf{score-based inference framework} that couples attribution-enhanced prompting, visual instruction scoring, and IQA metrics, providing an automatic way to support and evaluate cross-task VICL.
\end{itemize}

% In-context learning (ICL) enables foundation models, such as large language models (LLMs), to easily adapt to new tasks by leveraging a few input-output demonstrations without parameter updates~\cite{min2022rethinking, dong2022survey, highmore2024context, wang2023large, sia2024does, li2025large}, achieving strong zero-shot performance in similar linguistic contextual examples. 
% Recent progress in large vision-language models (VLMs) has extended this concept from language to multimodal (text and image) conditions, with works like MAE-VQGAN~\cite{bar2022visual}, Painter~\cite{wang2023painter}, Prompt Diffusion~\cite{wang2023incontext}, and X-Prompt~\cite{xprompt2024universal} showing great promise for visual in-context learning (VICL) from only a small set of visual demonstrations ~\cite{zhou2024visual, ma2024vision}.

% Intuitively, many visual problems share underlying relationships rather than existing in isolation ~\cite{zamir2018taskonomy, pal2019tmt, achille2019task2vec}. While fully supervised approaches typically learn each task independently, this paradigm is inefficient and demands large amounts of labeled data~\cite{lecun2015deep, chum2019beyond}. Task transfer learning addresses this challenge by exploiting knowledge from source tasks to accelerate and improve performance on new targets~\cite{bao2019information}. For example, the Taskonomy~\cite{zamir2018taskonomy} framework demonstrates that by computationally modeling transfer dependencies across diverse tasks, and the possibility to derive a taxonomic map that captures both direct and higher-order relationships has been proven. Such relational phenomena are hypothesized to exist in language domains as well. Viewed through the lens of language, many image-processing-oriented low-level tasks own the same prefix or suffix~\cite{trigka2025comprehensive, shu2024visual}, e.g. ``deraining'', ``denoising'', ``dehazing'', ``deblurring'', ``demoireing''. Textual descriptions of visual differences across such tasks reveal areas of overlap as well as points of divergence. In addition, textual descriptions of generation-oriented low-level tasks such as ``colorization'', ``light enhancement'', ``harmonization'', and ``style transfer'' ~\cite{bar2022visual, ke2023neural}, primarily hinge on linguistic variations in color, illumination, and contrast. This suggests that certain latent relationships may govern the transformations between them. Meantime, several studies~\cite{brooks2023instructpix2pix, potlapalli2023promptir, conde2024instructir} have demonstrated the regulatory role of language in guiding/driving visual expressions. 

% For a standard VICL inference framework (left half in Figure~\ref{fig:fig1}), a frozen vision or vision-language model accepts paired visual examples (input images and its labels) as the ``context'' along with a query image, and produces the output image for the query. Building on the above observations, we naturally raise the key question: \textit{If the visual prompt and the images to be learned come from two different tasks, can VLMs enable VICL across different tasks? }

% In this paper, we aim at providing a VLM-agnostic methodological analysis of generalization and task abstraction in cross-task scenarios, and developing an improved diagnostic framework to evaluate the capability of cross-task VICL in VLMs. Thus we propose T2T-VICL (right half in Figure~\ref{fig:fig1}), a collaborative pipeline based on multiple VLMs that enable VICL across multiple low-level cross-task pairs. We construct a novel approach that integrates implicit descriptive textural capture, gradient-based token attribution enhancement~\cite{sundararajan2017axiomatic, nguyen2025grains}, a visual instruction-guided metric~\cite{ku2023viescore}, and image quality assessment (IQA) metrics~\cite{zhang2012comprehensive}, allowing VLMs to perform cross-task in-context learning without the need for additional training and be evaluated. Our contributions can be summarized in three aspects: (1) We propose the first comprehensive multi-VLM pipeline that automatically generates implicit text prompts to distinguish two low-level vision tasks, which can be used to explore the \textbf{boundaries of cross-task VICL}. This work also introduces the \textbf{first text-image dataset} with cross-task implicit descriptions, establishing a benchmark expected to exert lasting influence within the field. (2) We design a \textbf{VLM $\rightarrow$ sVLM} and \textbf{sVLM $\rightarrow$ VLM} framework that gets knowledge from a large-scale model into a lightweight model and leverages the compact model to get the text prompt. (3) By coupling score-based reasoning, we establish an automatic inference framework that supports cross-task VICL in VLMs while \textbf{eliminating the need for further training}.
